% Source: 56_Chapter_43_Rootless_Immigrants_In_a_Sick_Society_A.txt
\bookchapter[Rootless Immigrants in a Sick Society]{CHAPTER 43}{ROOTLESS IMMIGRANTS IN A SICK SOCIETY}{American Evangelicals}

\leadin{D}{URING THE 1976 PRESIDENTIAL} campaign in the United States, Jimmy Carter’s evangelical
faith arose as one of the major issues. Aware that the Watergate scandal had toppled the Nixon
administration a short time before, many Americans felt that morality in government was of supreme
importance. Others, however, warned that religion could divide the nation and had no place in an
American political campaign.

Historian Arthur M. Schlesinger Jr., once an advisor to Presidents John Kennedy and Lyndon Johnson,
said bluntly, “I don’t think Carter should have brought in his intimacy with God, Jesus and so
forth~.~.~.~. If he feels that, bully for him, but it’s totally irrelevant to the campaign!”

In terms of a checklist for public service, Schlesinger was undoubtedly right. The American
Constitution prohibits religious tests for public office. But serious-minded Christians found it
impossible to look at government or any other social institution (the family, the courts, the
schools) through totally secular eyes.

That was the Evangelical’s dilemma in the Age of Ideologies. They offered millions of people a
consoling personal faith in the face of apocalyptic fears, but could they also lead America in
another great awakening of social ideals?

\section{RETURN OF THE OLD-TIME RELIGION}

In the 1970s, Jimmy Carter, the former Georgia governor, was a symbol of a revitalized evangelical
Christianity in the United States. The old-time religion was showing remarkable signs of new life.

Many Americans found this surprising. They had been brought up to think that revivalism was the
province of faith healers, holy rollers, and counterfeit preachers: a thing of high emotions, bad
taste, and simplemindedness. Yet in the 1970s prominent figures by the scores—in politics, sports,
and entertainment—spoke freely about their faith in Christ.

Religion that was supposedly confined to the Bible Belt suddenly appeared all over the country.
America, it seemed, had discovered born-again Christianity. After a generation of steady growth,
Evangelicals numbered forty-five million in the United States and had changed the face of American
Protestantism.

The so-called mainline—or to use Martin Marty’s term, public—Protestants belonged to the major
denominations in the National Council of Churches: Episcopalians, Methodists, Presbyterians,
American Baptists, United Church of Christ. But in most of the denominations Evangelicals formed a
conservative minority, so that about one-third of the thirty-six million Protestants in the National
Council were evangelical. When these were added to about thirty-three million Evangelicals outside
the Council, the so-called Private Protestants constituted a major religious body, close in number
to the forty-six million Roman Catholics.

Of course, Evangelicals represented no single approach to America’s problems. The movement included
a number of distinct subgroups. Among them were the fundamentalists, the militant right wing
churchmen who opposed all accommodation to contemporary culture, and the Pentecostals, who had
experienced the “baptism in the Holy Spirit” and practiced such divine gifts as speaking in tongues
and miraculous healing by prayer. These Pentecostals, or charismatics as some preferred, included
everything from Episcopalians and nearly a million Roman Catholics, to faith healers and assorted
tent preachers. Most Evangelicals, though, were basically conventional Protestants who held
staunchly to the authority of the Bible and to orthodox Christian doctrines. They believed in making
a conscious personal commitment to Christ: a spiritual encounter, gradual or instantaneous, known as
the born-again experience.

\begin{sidebar}{Who Are the Charismatics?}
The term charismatic is used with great variety. There are charismatic denominations such as
Assemblies of God, but charismatics also make up significant percentages of other denominations,
such as the Episcopalians.

As a simple way to begin, people typically are described as charismatic if the exercise of the more
dramatic spiritual gifts, such as tongue speaking, healing, and acts of prophecy, is an important
and routine part of the church’s worship and an individual believer’s personal devotion.
\end{sidebar}

In spite of the growth of the evangelical movement and its sudden visibility in the media it gave
little evidence of reshaping the ideas and ideals of American culture. It had little impact in the
government or the universities or the communication centers of the nation. The reasons for
estrangement from politics, education, and the media lie in the early years of the twentieth
century.

Historian George M. Marsden of Notre Dame compares Evangelicals in twentieth-century America to
rootless immigrants in a new land. Only in this case Evangelicals never crossed the ocean to
experience the shock of a new land. They simply held their ground while the rest of the country
changed.

As we have seen, after the Civil War Southern Evangelicalism was battered by defeat and a sense of
hopelessness. Much of the northern wing turned to premillennialism, the belief that Christ’s return
was imminent and that society would inevitably get worse before it occurred. By the late 1800s the
great evangelist Dwight L. Moody literally preached a lifeboat ethic: “I look on the world as a
wrecked vessel. God has given me a lifeboat and said, ‘Save all you can.’\,” Thus many conservatives
withdrew from the social arena. Evangelical historian Timothy Smith describes this as the Great
Reversal.

The roots of fundamentalism lie in this surrender of social concern. In the 1880s and 1890s Bible
study and personal holiness seemed more rewarding than the reform of American life.

Beginning with a small meeting at Swampscott, Massachusetts, in July 1876, Evangelicals from many
denominations gathered in summer Bible conferences to discuss the second coming of Christ. The
Baptist minister A. J. Gordon of Boston played a leading role in organizing two large prophecy
conferences in New York (1878) and Chicago (1886) at which the essentials of premillennialism were
hammered out: The world will continue to decline into sin until the Antichrist is loosed for one
last orgy of destruction, at the end of which time Christ will return with his saints to establish
an earthly reign of one thousand years. A peculiar form of premillennialism was articulated with
growing clarity in such conferences. Its name, dispensationalism, is not very suggestive, but the
last three dispensations (ages) before eternity are crucial to its classical version. Jesus offered
himself as a literal king for the nation of Israel to rule from Jerusalem; the Hebrew people
rejected his offer. This rejection marked the end (or the postponement) of the first of these three
ages. The next age was the age of the church; being rejected as Israel’s king, Jesus turned his
attention to the non-Hebrews. To these Gentiles he offered another covenant of grace and forgiveness
that followed the covenant of law and obedience. Eventually, this age runs its course, ending in
failure, as does every covenant. Jesus’ second return marks the beginning of the last of the three
ages: he returns as king and rules through Israel over the entire world for one thousand years. This
version of premillennial thinking was a serious departure from historic premillennialism; it
demanded that God had two different peoples, Israel and the church. These two peoples stem from two
completely unrelated covenants, which operated by completely different economies: grace versus
obedience to the law. In its original form, it breaks several longstanding beliefs: that Jesus and
his church fulfilled the Old Testament promises, the unity of the Bible, grace in the Old Testament,
and Jesus’ call for obedience in the New Testament. Additionally, there was much in the New
Testament that was for Israel and not for the church, most remarkably the Sermon on the Mount. Many
amendments to this theory have made it more acceptable to more orthodox-minded believers (such as
Presbyterians). Today it is widely held among the conservative Protestants. Ironically, it shares
some convictions with modern thinkers: with liberals, it holds that Jesus was initially to be
understood in earthy, political terms; it also explains (in concert with other ideas, such as the
cessation of gifts) why the current church life seemed so different from the experience reflected in
the Bible.

\tabletitle{Three Crucial Ages of Dispensationalism}
\begin{booktable}{@{}L{0.16}L{0.24}L{0.24}L{0.24}@{}}
\toprule
 & .~.~. Jewish age & Church age & 1,000 year rule \\
\midrule
People & Israel & Church & Israel/ World \\
\addlinespace[2pt]
Predominant Race & Hebrews & Gentiles & Hebrews/ World \\
\addlinespace[2pt]
Jesus’ role & Jesus comes to be king over Israel. & After rejection, Jesus turns to Gentiles and mercifully dies to pay for sin. & Jesus returns to rule for 1,000 years. \\
\bottomrule
\end{booktable}

During the same years other Evangelicals turned to holiness conferences. Concern for a second
blessing, entire sanctification, or Christian perfection had always been a main tenet of Methodist
revivalism. In the late nineteenth century members of other religious communions came to share this
concern. Holiness groups, such as the Church of the Nazarene, as well as deeper life conferences,
urged believers to yield themselves to reliance on the Holy Spirit as a way to find the victorious
Christian life.

While increased attention both to the end times and to personal Christian living had firm biblical
roots, it also gave traditional Evangelicals a way of maintaining their faith in a culture over
which they were steadily losing control. If they could not shape the affairs of men, they could find
comfort in the world of the Spirit.

\section{THE APPEARANCE OF FUNDAMENTALISM}

Fundamentalism is usually dated from a series of twelve small books published from 1910 to 1915
containing articles and essays designed to defend fundamental Christian truths. Three million copies
of the books were sent free to theological students, Christian ministers, and missionaries all over
the world.

The project arose in the thinking of Lyman Stewart, a wealthy oilman in Southern California, who was
convinced that something was needed to reaffirm Christian truths in the face of biblical criticism
and liberal theology. After listening to the Reverend Amzi C. Dixon preach in August 1909, Stewart
secured Dixon’s help in publishing The Fundamentals.

Stewart then enlisted the financial support of his brother, Milton, and Dixon, pastor of the Moody
Church in Chicago, chose a committee, which included the evangelist R. A. Torrey, to assist in the
editorial work.

Sixty-four authors were eventually chosen to appear in The Fundamentals. The American premillennial
movement and the English Keswick Conference were well-represented. Other conservatives, however,
were also among the contributors, including E. Y. Mullins of Southern Baptist Seminary and Benjamin
B. Warfield of Princeton Seminary.

World War I delayed the outbreak of the “modernist-fundamentalist” controversy in the Protestant
denominations. But shortly after the doughboys came home from Europe, Baptists, Presbyterians,
Methodists, and the Disciples of Christ launched their own war of words over the values and the
dangers of liberal theology in the churches.

In 1920, Curtis Lee Laws, Baptist editor of the Watchman-Examiner, called “fundamentalists” within
the Northern Baptist Convention to a conference in Buffalo, New York. This group of conservatives,
popularly called The Fundamentalist Fellowship, were moderate conservatives. They believed that the
modernists were surrendering the fundamentals of the gospel: the sinful nature of man, his inability
to be saved apart from God’s grace, the centrality of Jesus’ death for the regeneration of the
individual and the renewal of society, and the authoritative revelation of the Bible. This group was
the first to apply the name fundamentalist to itself. They were unable to gain the adoption of a
confession of faith among Northern Baptists, but Laws and his associates did not consider their
cause a lost one. In 1924 Laws wrote that certain schools of his denomination had checked the
inroads of liberalism and that the investigation of the mission societies, advocated by the
fundamentalists, resulted in changes that made the creation of a new mission unnecessary. More
militant Baptists disagreed with Laws and formed the fundamentalist General Association of Regular
Baptists for fundamentalists.

The Presbyterian champion of orthodoxy was Professor J. Gresham Machen of Princeton Theological
Seminary. In 1929 the General Assembly of the Presbyterian Church authorized a reorganization of the
seminary. Machen and a small retinue of distinguished professors at the school felt that a merger of
boards strengthened the liberal influence in the school. They withdrew from Princeton in protest and
founded Westminster Theological Seminary in Philadelphia.

When Machen refused to break his ties with the Independent Board of Presbyterian Foreign Missions,
he was brought to trial in his church’s courts and found guilty of rebellion against superiors. As a
result conservatives in the denomination founded the Orthodox Presbyterian and the Bible
Presbyterian churches.

At the heart of these modernist-fundamentalist differences were conflicting views of the Bible.
While a variety of positions could be found in each camp, Shailer Mathews and Machen may be taken as
representatives of the two parties that emerged.

Mathews was professor of historical theology and Dean of the Divinity School at the University of
Chicago. In 1924 his book The Faith of Modernism set forth his case for a view of the Bible based
upon “scientific investigations.” Confessional theology, he said, considers Scripture supernaturally
given.

\begin{blockquote}
The modernist uses Scripture as the trustworthy record and product of a developing religion~.~.~.~.
In discovering this experience of God and accepting it as his own religious ancestry the modernist
affirms the trustworthiness of the Scripture~.~.~.~. Christianity becomes not the acceptance of a
literature but a reproduction of attitudes and faith, a fellowship with those ancient men of
imperfect morals whose hearts found God.
\end{blockquote}

All the hallmarks of liberal theology are here: (1) the evolutionary philosophy applied to religion;
(2) the optimistic view of man centering in his religious experience; and (3) the moralistic
conception of God, who can so readily be “found” by man.

In 1915, in The Princeton Theological Review, Professor Machen, probably the most articulate
defender of orthodoxy, answered the modernist’s appeal to historical and literary methods. “The
student of the New Testament,” he said, “should be primarily an historian.~.~.~. The Bible contains
a record of something that has happened, something that puts a new face upon life~.~.~. It is the
life and death and resurrection of Jesus Christ. The authority of the Bible should be tested here at
the central point. Is the Bible right about Jesus?~.~.~. A teacher and example, or a Savior?”

Here, too, are the marks of early fundamentalism: (1) a supernatural Jesus attested by his
resurrection from the dead; (2) a trustworthy Bible, the fountain of the Christian faith; and (3)
the need of men to have “a new face upon life.”

\section{WILLIAM JENNINGS BRYAN, SYMBOL OF CHANGE}

Squalls in the denominations, however, were indicators of wider winds of change in America. If we
looked for one man whose life illustrates the shift in American views, few would serve as admirably
as William Jennings Bryan.

Thanks to television, Broadway, and the movie Inherit the Wind, the mere mention of Bryan’s name
recalls that packed courtroom in Dayton, Tennessee, where, on a hot July day in 1925, Bryan and
Clarence Darrow nearly came to blows over the Genesis account of man’s origin.

The setting was the trial of John Scopes, who had been arraigned for violating Tennessee’s law
against the teaching of evolution in the public schools.

Darrow, at the height of his powers as a courtroom artist, confused and embarrassed the aging
Bryan’s attempts to defend the biblical record. A press corps overwhelmingly hostile to Bryan’s
point of view disgorged millions of words of coverage to a watching nation. The image of a
sweat-streaked Bryan wilting before the searing rational assault from Darrow maintains its overrated
and largely undeserved place as a symbol of the stupidity of Bible-believing American Christians.
Ironically, it was Bryan who was the true political progressive. Inherit the Wind leaves a
conspicuously deceptive impression. One would gather that Bryan was a hyper-conservative naysayer
and Darrow was the progressively minded forward-thinker. Writers invented a story line that Bryan
had bullied Mr. Scopes’s fiancée to legitimize the trial strategy of demonizing Bryan.

The irony of Bryan’s image as a bigoted fanatic arises from the fact that all of his life he stood
for the common man. No one fought harder for “the average American.” It was Bryan who was against
social Darwinism and championed women’s right to vote. These two ideas are linked: if emotionally
oriented women had not evolved enough to be sufficiently rational then they should not be allowed to
vote. As leader of the Progressive cause in the Democratic party, Bryan was a three-time candidate
for president of the United States, and secretary of state in the cabinet of President Woodrow
Wilson. As Wilson’s secretary of state Bryan expressed his devotion to peace by negotiating
arbitration treaties with thirty countries. When the crisis of 1914 broke, however, his treaties
proved ineffective after the Germans sank the Lusitania on May 7, 1915. President Wilson and the
cabinet opposed the use of the treaties. This conflict of policy led to Bryan’s resignation.

This near end of his political career, however, only meant new opportunities in Bryan’s reforming
and religious career. He soon threw himself into the Prohibition cause and played no small role in
securing passage of the Eighteenth Amendment outlawing alcoholic beverages across the country after
January 1920. This was probably the last successful evangelical crusade for a moral America.

Later, Americans tended to view Prohibition as the austere weapon of a pack of bluenosed Puritans
who found life a joyless thing and were determined that no one else should be allowed to distill a
drop of pleasure from it. That attitude is itself an example of the shift in American culture.

Bryan saw Prohibition as another struggle against those selfish interests that put private profit
above human welfare and feed upon the helplessness of the masses. Those who labor for Prohibition,
he wrote, “are helping to create conditions which will bring the highest good to the greatest
number,~.~.~. for it is not injustice to any man to refuse him permission to enrich himself by
injuring his fellowmen.”

Bryan’s last great crusade was the one that drew him directly into the fundamentalist movement, the
effort to outlaw the teaching of evolution in the public schools of America.

Prior to the 1920s Bryan frequently revealed his awareness of the teachings of Darwin and the
dangers they posed for public morality. As early as 1905, after reading Darwin’s Descent of Man, he
said that the biologist’s conception of man’s origin would weaken the cause of democracy.

In the spring of 1921 Bryan issued a series of attacks on evolution that instantly placed him in the
forefront of fundamentalist forces. Whether evolution was right or wrong was not his concern. He
never intended to refute evolution on scientific grounds. “The objection to evolution,” he said, “is
not, primarily, that it is not true. The principle objection is that it is highly harmful to those
who accept it.”

His lecture, “The Menace of Darwinism,” was printed in pamphlet form and widely distributed.
Morality and virtue, he wrote, are dependent upon religion and a belief in God, and therefore
anything that weakens belief in God weakens man and makes him unable to do good. By putting man on
the level of a brute and ignoring all spiritual values, evolutionary theory applied to social theory
robbed man of his major stimulus to moral living. Without a moral compass it seemed reasonable that
the strong should flourish and not be inhibited by concern for the weak.

Legal opposition to evolution took the form of thirty-seven anti-evolution bills introduced into
twenty state legislatures between 1921 and 1929. They were drawn up by a generation that ratified
the Eighteenth Amendment, products of an American belief that legislative action can assist public
morality.

The Scopes Trial was the consequence of the law adopted in Tennessee, which made it unlawful for a
public school teacher “to teach the theory which denies the story of divine creation of man as
taught in the Bible, and to teach instead that man has descended from a lower order of animals.”

The immediate issue before the Dayton court was whether John Scopes, the young high school biology
teacher, had violated the Tennessee law. But the legal question was strictly secondary. Businessmen
and reporters turned the trial into a circus. The star performers were Bryan, who served as a
prosecuting attorney, and Clarence Darrow, Chicago’s brilliant, well-known lawyer, who defended
Scopes. Both men saw the larger dimensions of the trial.

“My only purpose in coming to Dayton,” said Bryan with a glance at Darrow, was “to protect the Word
of God against the greatest atheist or agnostic in the United States!” The courtroom broke into
applause. Darrow argued that nothing less than intellectual freedom was on trial and proceeded to
make Scopes the real plaintiff and Bryan the real defendant. He used Bryan’s own testimony as
evidence of fundamentalist stupidity.

When all was over, Scopes was found guilty and fined a token sum. Bryan, thus, won in Dayton, Darrow
in the rest of the country. Five days after the trial Bryan passed away quietly in his sleep,
leaving all of his reform causes behind. In a very real sense the evangelical crusade for a
Christian America died with him.

A short time later the Great Depression hit America hard and evangelical Christianity struggled with
the rest of the country. It survived in assorted voluntary societies and radio broadcasts, but the
public often stereotyped fundamentalists as closed-minded, ignorant, belligerent, and separatistic.

\section{THE RENEWAL OF EVANGELICAL CHRISTIANITY}

After World War II, however, evangelical Christianity returned to public prominence in America. Its
best known voice had a soft North Carolina accent. Billy Graham became a household name by preaching
to thousands in every major stadium in the United States, by a regularly scheduled radio broadcast,
and by telecasts on national networks. Decision magazine, published by his evangelistic association,
went to multiplied millions of homes.

Graham had grown up in the Southern Baptist and Southern Presbyterian tradition. He had trained at
Bob Jones University, Florida Bible Institute, and Wheaton College. After a pastorate in Western
Springs, Illinois, a suburb of Chicago, he entered full-time evangelistic work with a youth
organization called Youth for Christ. After a highly successful and widely publicized evangelistic
campaign in Los Angeles (1949) Graham’s fame spread rapidly. He was soon conducting citywide
crusades around the world.

Billy Graham, however, was only the most prominent name in a new style evangelicalism. Many
conservative Christians expressed dissatisfaction with the drift of fundamentalism between the two
world wars. Most had no criticism of the doctrines of the movement, but they felt it had stumbled
into the mire of some unfortunate emphases.

As early as 1947 Carl F. H. Henry in his Uneasy Conscience of Modern Fundamentalism expressed
distress over fundamentalism’s failure to apply basic Christian truths to crucial problems
confronting modern man.

In 1956 Billy Graham led a group of Evangelicals in launching a new periodical called Christianity
Today. Carl Henry left his teaching post at Fuller Theological Seminary to serve as editor. In the
beginning the magazine was sent without charge to theological students and clergymen of various
doctrinal persuasions, assuring it of a wide circulation.

In the first issue Henry said the magazine aimed to express historical Christianity to the present
generation. The founders felt that theological liberals had failed to meet the moral and spiritual
needs of people. The editors accepted unreservedly the “complete reliability of the written Word of
God.” But they intended to present the implications of the gospel for every area of life.

The evangelical awakening gained unexpected reinforcement in 1960 from another source of personal
Christianity, Pentecostalism. The Pentecostal experience—“the baptism of the Holy Spirit with the
evidence of speaking in tongues”—was not new. The spark of twentieth-century Pentecostalism was a
three-year-long revival, beginning in 1906, at the Azusa Street Mission in Los Angeles. There were
personal experiences of tongue speaking earlier, but Azusa Street ignited worldwide Pentecostalism.
Christians from all over North America, Europe, and the Third World visited Azusa Street and carried
the fire back home.

Pentecostal denominations followed. The largest in the United States were the Assemblies of God, the
Church of God in Christ, the Church of God, and the Pentecostal Holiness Church. These were usually
filled with socially and economically depressed people. That is the reason the secular press took
note in 1960 when the Pentecostal experience broke out in middle-class Lutheran and Episcopal
congregations. That made hot copy.

This new explosion of Pentecostal passion was dubbed neo-Pentecostalism. Its fallout drifted in all
directions. Soon America was dotted with charismatic prayer groups, composed mostly of members of
mainline Protestant churches. They claimed to have found a new dimension in their Christian lives
through the Pentecostal baptism.

Among the Lutherans, the American Lutheran Church saw most of the action in the early 1960s.
Controversy that erupted in 1963 was only subdued when authorities issued guidelines for practicing
Pentecostals. Then in 1967 the experience touched Missouri Synod Lutheran circles. Several of its
ministers were defrocked because of public statements on the issue.

The extent of Pentecostal penetration of Lutheran churches was suggested in 1972 by the more than
ten thousand who attended the First International Lutheran Conference on the Holy Spirit in
Minneapolis, Minnesota.

In 1970 the General Assembly of the United Presbyterian Church in the U.S.A. issued a report on “The
Work of the Holy Spirit.” Reflecting two years of careful study, it was hailed as one of the best
official statements on the charismatic renewal and brought about considerable growth in the
movement. By 1975 some spokesmen estimated that between ten thousand and fifteen thousand members of
the United Presbyterian Church and the Presbyterian Church in the U.S.A. had received the
Pentecostal baptism. These two denominations—Lutherans and Presbyterians—are only representative of
the growth in other mainline Protestant denominations and in Roman Catholicism.

As the 1970s closed in the United States, Evangelicals of all varieties were troubled by the
question, does faith matter? For decades they had assumed that significant numbers of born-again
Christians would make a difference in the moral climate of America. The 1970s, however, brought that
assumption to a severe test. Evangelicals had their numbers, but American society was plagued by
crime, divorce, racism, violence, sexual perversion, alcoholism, and drugs. Some observers wondered
if Arthur Schlesinger Jr. could have been speaking not only of Jimmy Carter but of all Evangelicals:
“If you feel that, bully for you, but it is totally irrelevant.”

\readinghead{Suggestions for Further Reading}
\begin{readinglist}
\entry Levine, Lawrence W. Defender of the Faith. New York: Oxford, 1965.
\entry Marsden, George M. Fundamentalism and American Culture. New York: Oxford, 1980.
\entry Noll, Mark, David W. Bebbington, and George A. Rawlyk, eds. Evangelicalism. New York: Oxford,
  1994.
\entry Quebedeaux, Richard. The New Charismatics. Garden City, NY: Doubleday, 1976.
\entry Sandeen, Ernest R. The Roots of Fundamentalism. Chicago: University of Chicago Press, 1970.
\entry *Synan, Vinson. The Century of the Holy Spirit: 100 years of Pentecostal and Charismatic
  Renewal, 1901–2001. Nashville: Thomas Nelson, 2001.
\entry Woodbridge, John, Mark A. Noll, and Nathan O. Hatch. The Gospel in America. Grand Rapids:
  Zondervan, 1979.
\end{readinglist}
